\section{Gaussian 图模型用精度矩阵编码条件独立}
\label{sec:13-Machine-Learning/07-Graphical-and-Sequential-Models/02}

你拿到六十个脑区的 fMRI 时间序列，想画一张``哪些脑区直接耦合''的网络图。最顺手的做法是算相关矩阵，把相关系数超过阈值的那些对连起来。结果是一张几乎全连的图：\(1770\) 对里有一千多对显著。降低阈值图更糊，提高阈值又只剩下几条最强的边，中间那些说不清是直接耦合还是隔着第三个脑区传过来的。相关系数回答的是两个变量之间有没有关联，而你想问的是扣掉所有别的脑区之后还剩不剩关联——这是两个不同的问题。

上一节说，无向图的边表示给定其余全部变量后仍然存在的依赖。在联合 Gauss 的假设下，这句话可以落到矩阵元素上：协方差的零对应边缘独立，精度矩阵的零对应条件独立。

\subsection{同一个分布，协方差稠密而精度稀疏}

设 \(X\sim\mathcal N(\mu,\Sigma)\)，记 \(K=\Sigma^{-1}\) 为精度矩阵。两种零元素承担不同语义：

\[
\Sigma_{ij}=0
\quad\Longleftrightarrow\quad
X_i\perp X_j,
\qquad
K_{ij}=0
\quad\Longleftrightarrow\quad
X_i\perp X_j\given X_{-\set{i,j}}.
\]

左边那条依赖 Gauss 分布里``不相关即独立''这一特权，换成别的分布就不成立；右边那条才是无向图缺边所用的关系。

差别有多大，看一个四变量的例子。取平稳 AR(1)：\(X_t=\rho X_{t-1}+\varepsilon_t\)，\(\rho=0.8\)。相关矩阵的元素是 \(\rho^{\abs{i-j}}\)，于是 \(\operatorname{Corr}(X_1,X_4)=0.512\)——放进相关网络里，这是一条相当粗的边。而 \(K\) 是三对角的，\(K_{14}=0\)，意味着给定 \(X_2,X_3\) 之后 \(X_1\) 与 \(X_4\) 条件独立。那条 \(0.512\) 完全是沿着链传过来的。

\begin{center}
\begin{tikzpicture}[line width=0.7pt]
\foreach \i/\j/\s in {0/0/72,1/1/72,2/2/72,3/3/72,
  0/1/56,1/0/56,1/2/56,2/1/56,2/3/56,3/2/56,
  0/2/40,2/0/40,1/3/40,3/1/40,
  0/3/26,3/0/26}
  {\fill[black!\s] (0.5*\j,-0.5*\i) rectangle ++(0.5,-0.5);}
\draw[step=0.5,black!55] (0,-2) grid (2,0);
\node at (1,0.4) {\(\Sigma\)：十六个元素全非零};
\begin{scope}[shift={(4.4,0)}]
  \foreach \i/\j/\s in {0/0/72,1/1/72,2/2/72,3/3/72,
    0/1/56,1/0/56,1/2/56,2/1/56,2/3/56,3/2/56}
    {\fill[black!\s] (0.5*\j,-0.5*\i) rectangle ++(0.5,-0.5);}
  \draw[step=0.5,black!55] (0,-2) grid (2,0);
  \node at (1,0.4) {\(K=\Sigma^{-1}\)：只剩三对角};
\end{scope}
% 左：相关网络，四点全连
\begin{scope}[shift={(0,-3.3)}]
  \draw (0,0) -- (0.667,0);
  \draw (0.667,0) -- (1.333,0);
  \draw (1.333,0) -- (2,0);
  \draw (0,0) .. controls (0.667,0.62) .. (1.333,0);
  \draw (0.667,0) .. controls (1.333,0.62) .. (2,0);
  \draw (0,0) .. controls (1,-0.85) .. (2,0);
  \foreach \x/\lab in {0/1,0.667/2,1.333/3,2/4}{
    \fill[white] (\x,0) circle (0.18);
    \draw (\x,0) circle (0.18);
    \node[font=\small] at (\x,0) {\lab};
  }
  \node at (1,-1.4) {按相关连边};
\end{scope}
% 右：条件依赖网络，链
\begin{scope}[shift={(4.4,-3.3)}]
  \draw (0,0) -- (0.667,0);
  \draw (0.667,0) -- (1.333,0);
  \draw (1.333,0) -- (2,0);
  \foreach \x/\lab in {0/1,0.667/2,1.333/3,2/4}{
    \fill[white] (\x,0) circle (0.18);
    \draw (\x,0) circle (0.18);
    \node[font=\small] at (\x,0) {\lab};
  }
  \node at (1,-1.4) {按精度连边};
\end{scope}
\draw[black!35] (3.2,-4.9) -- (3.2,0.7);
\end{tikzpicture}
\end{center}

\emph{同一个 AR(1) 分布：相关矩阵处处非零，精度矩阵只留下相邻项。}

图里左右两侧描述的是同一个四维 Gauss 分布，只是问了不同的问题。左边那张全连图没有说错任何事实，它如实报告了每一对都相关；它只是没有回答你关心的那个问题。

\subsection{条件化就是取 Schur 补}

把变量分成两块 \(X_A,X_B\)，协方差相应分块。若 \(\Sigma_{BB}\) 可逆，则

\[
X_A\given X_B=x_B
\sim
\mathcal N\bigl(
\mu_A+\Sigma_{AB}\Sigma_{BB}^{-1}(x_B-\mu_B),
\;
\Sigma_{AA}-\Sigma_{AB}\Sigma_{BB}^{-1}\Sigma_{BA}
\bigr).
\]

条件协方差是一个 Schur 补。观察 \(X_B\) 只削掉能被 \(B\) 线性解释的那部分不确定性，剩下的方差不会因为``已经看过很多变量''而自动趋于零。

取 \(A=\set{i}\)、\(B\) 为其余全部变量，用精度形式整理同一个式子，得到

\[
X_i\given X_{-i}=x_{-i}
\sim
\mathcal N\bigl(
\mu_i-K_{ii}^{-1}K_{i,-i}(x_{-i}-\mu_{-i}),
\;
K_{ii}^{-1}
\bigr).
\]

条件均值里只出现 \(K_{ij}\neq 0\) 的那些变量。节点 \(i\) 的邻居就是精度矩阵第 \(i\) 行的非零非对角位置，前一段的三对角结构因此直接读成一条链。与之配套的是偏相关系数

\[
\rho_{ij\cdot-\set{i,j}}
=-\frac{K_{ij}}{\sqrt{K_{ii}K_{jj}}},
\]

它为零当且仅当 \(K_{ij}=0\)。这条等价随 Gauss 假设成立而成立：离开 Gauss 族，偏相关为零只排除了线性的条件关系，非线性依赖照样可以存在。

\subsection{二次型展开后，矩阵元素就是图上的势}

中心化之后 Gauss 密度写成 \(p(x)\propto\exp(-\tfrac12 x^{\trans}Kx)\)，把二次型展开：

\[
-\frac12x^{\trans}Kx
=-\frac12\sum_i K_{ii}x_i^2
-\sum_{i<j}K_{ij}x_ix_j.
\]

对角项是单节点势，非零的 \(K_{ij}\) 是一条边上的成对相互作用。若 \(K_{ij}=0\)，密度里根本没有 \(x_ix_j\) 这一项，图上也就没有这条边——上一节讲的无向图分解，在这里退化成了矩阵元素的读法。

约束 \(K\succ0\) 同时保证积分收敛与协方差存在，也意味着不能随手写一个稀疏对称矩阵当精度用：稀疏容易，正定不容易，两者要一起满足。

Gauss 情形与一般无向图模型的分水岭在归一化常数上。上一节说配分函数 \(Z\) 要对全空间求和，在这里它有闭式：

\[
A(K,\mu)=\frac d2\log(2\pi)-\frac12\log\det K+\frac12\mu^{\trans}K\mu .
\]

代价从指数级枚举降到一次 Cholesky 分解，量级是 \(O(d^3)\) 甚至更少（矩阵稀疏时）。连续、二次、正定这三个条件换来了可算的 \(Z\)，这也是 Gauss 图模型能在几万个节点上做精确推断、而离散 Markov 随机场往往做不到的原因。

\subsection{有向 Gauss 模型道德化后才是无向图}

给变量选定拓扑顺序，可以写线性结构方程

\[
X_i=\sum_{j\in\operatorname{pa}(i)}B_{ij}X_j+\varepsilon_i,
\qquad
\varepsilon_i\ \text{相互独立},\ \varepsilon_i\sim\mathcal N(0,\omega_i^2).
\]

向量形式为 \(X=BX+\varepsilon\)，按拓扑顺序排列后 \(B\) 严格下三角、\(\Omega=\Var(\varepsilon)\) 对角。解出 \(X=(I-B)^{-1}\varepsilon\)，于是

\[
\Sigma=(I-B)^{-1}\Omega(I-B)^{-\trans},
\qquad
K=(I-B)^{\trans}\Omega^{-1}(I-B).
\]

第二个式子值得盯一会儿。\(K_{ij}\) 是 \((I-B)\) 的第 \(i\) 列与第 \(j\) 列经 \(\Omega^{-1}\) 加权的内积，只要存在某个 \(k\) 同时以 \(i\) 和 \(j\) 为父节点，这个内积就一般不为零。换句话说，DAG 里两个没有连边的父节点，在无向图里会被连上——这正是上一节碰撞结构的矩阵版本：给定共同子节点会打开通道，而无向图的边描述的恰好是给定其余全部变量后的依赖。把 DAG 转成无向图时先给每个节点的父节点两两连边、再抹去方向，这一步叫道德化，它损失的信息就是碰撞结构。

\subsection{高维下样本协方差不可逆，稀疏得靠正则化}

给定中心化样本的经验协方差 \(S\)，对数似然中与精度有关的部分为

\[
\ell(K)\propto\log\det K-\trace(SK),\qquad K\succ0,
\]

在 \(S\) 可逆时的最大值点是 \(K=S^{-1}\)。可是 \(d\) 接近或超过样本量 \(n\) 时 \(S\) 奇异，逆根本不存在；即便勉强可逆，最小特征值附近的噪声也会被求逆放大，产生成片的虚假强边。六十个脑区、两百个时间点已经落在这个尴尬区间里。

图 lasso 在似然上加一个非对角的 \(L_1\) 惩罚：

\[
\widehat K\in\argmax_{K\succ0}
\left[\log\det K-\trace(SK)-\lambda\sum_{i\ne j}\abs{K_{ij}}\right].
\]

惩罚项把小的 \(K_{ij}\) 压成精确的零，于是解本身就带着一张图。\(\lambda\) 越大图越稀疏，但它是偏差与方差之间的一个选择，不是``真实边数''的客观刻度：换一个 \(\lambda\)，边集就换一批。惩罚还对变量尺度敏感，通常需要先标准化，否则单位大的变量会被系统性地判为无边。

计算上不要先求 \(S^{-1}\) 再修补。Cholesky 分解、对数行列式与线性方程求解都能直接利用正定结构，数值上稳得多。拿到解之后值得再看三个数：最小特征值离零有多远，条件数有多大，以及在自助重抽样下同一条边出现的频率。

\subsection{一张稀疏图是估计结果，不是机制图}

前面每一步都挂着条件。联合 Gauss 若不成立，精度的零不再等价于条件独立，只等价于偏相关为零；重尾或非线性依赖会让两者分家。若存在未观测的共同原因，它会在观测变量的精度矩阵里留下一片非零，逼算法补上现实中不存在的直接连接。测量误差和样本选择也会改变边集。

因此报告一张 Gauss 图模型时，四层条件最好分开写清楚：联合 Gauss 是否站得住，精度用什么方法估计，\(\lambda\) 依据什么选定，边在重抽样或新数据上是否复现。边的方向更是另一回事——无向图本来就不带方向，而结构方程里的方向来自你选的拓扑顺序，观察数据无法把它定下来。

链结构在本节只是个玩具例子，实际用途却不小：只要把节点排成时间轴，三对角精度矩阵描述的就是一条线性动态系统的轨迹。下一节让这条链动起来，看看在时间方向上做 Gauss 条件化会得到什么。
